%% ============================================================
%%  CHAPTER 9 — CLASSIFICATION AND TRAINING
%% ============================================================
\chapter{Classification and Training}
\label{ch:classification}

This chapter presents the classification and training methodology
used to learn the mapping from the engineered feature superset
(Chapter~8) --- and, for two of the five gesture classes, from the
raw preprocessed waveform directly --- to the five final gesture
classes established in Chapter~6. This is the chapter in which the
project's central engineering narrative plays out: an initial flat
classifier is shown to be evaluated dishonestly by a naive
cross-validation protocol; a rigorous leakage audit is performed to
establish what the model can actually be trusted to do; and the
final architecture is built not by adding more features to one
classifier, but by routing each gesture family to the
classification method that is physiologically matched to it.
Section~\ref{sec:cl_overview} provides an overview of the training
pipeline and its evolution across four production scripts.
Section~\ref{sec:cl_imbalance} describes class imbalance handling
via \acrshort{smote}. Section~\ref{sec:cl_leakage} presents the
evaluation methodology crisis and the four-test leakage audit that
resolved it --- a central methodological contribution of this
project. Section~\ref{sec:cl_loso} defines the
Leave-One-Session-Out (\acrshort{loso}) protocol used for all
honest reporting from this point forward.
Section~\ref{sec:cl_baseline} presents the flat Random Forest
baseline. Section~\ref{sec:cl_hierarchical} presents the
hierarchical, physiology-routed architecture and its reconciliation
under \acrshort{loso}. Section~\ref{sec:cl_riemann} derives the
Riemannian-geometry muscle sub-classifier. Section~\ref{sec:cl_dtw}
tells the five-attempt story of the blink sub-classification
ceiling and its resolution via \acrshort{dtw} shape-matching.
Section~\ref{sec:cl_architecture} assembles the final three-way
hybrid architecture. Section~\ref{sec:cl_progression} presents the
cumulative performance progression. Section~\ref{sec:cl_summary}
summarises the chapter.

%% ----------------------------------------------------------
\section{Training Pipeline Overview}
\label{sec:cl_overview}
%% ----------------------------------------------------------

The training codebase evolved through four production scripts,
each corresponding to a stage of the architectural narrative told
in this chapter:

\begin{itemize}[leftmargin=*]
    \item \code{train.py} --- trains the flat Random Forest
    baseline on the full feature superset (Section~\ref{sec:cl_baseline}).
    \item \code{train\_hierarchical.py} --- trains the two-stage
    hierarchical router (Section~\ref{sec:cl_hierarchical}).
    \item \code{train\_hybrid.py} --- replaces the muscle
    sub-classifier's hand-crafted features with Riemannian
    covariance geometry (Section~\ref{sec:cl_riemann}).
    \item \code{train\_dtw\_hybrid.py} --- replaces the blink
    sub-classifier with \acrshort{dtw} shape-matching, producing the
    final production model (Section~\ref{sec:cl_dtw}).
\end{itemize}

All four scripts share the same upstream data pipeline: the
activity-gated epochs and 129-feature superset from Chapters~6--8
are loaded once, and each script fits \code{StandardScaler},
class-balancing, and the relevant classifier(s) strictly on
whichever training fold is active for the evaluation protocol in
use (Section~\ref{sec:cl_loso}). \acrshort{smote} is applied to the
training fold only, after any split, exactly as in the original
flat pipeline (Section~\ref{sec:cl_imbalance}).

Two evaluation protocols recur throughout this chapter and should
not be confused: \textbf{GroupKFold} (recording-level, no epoch
overlap leakage) is used during development and ablation studies,
where the goal is fast iteration; \textbf{Leave-One-Session-Out}
(\acrshort{loso}) is used for every number reported as a final,
honest result. The reasoning behind this choice, and the audit
that led to it, is the subject of
Sections~\ref{sec:cl_leakage}--\ref{sec:cl_loso}. A separate,
free-running mixed-gesture validation recording was also attempted
(\code{Random/}, Chapter~6) but could not be used for evaluation
due to an unsynchronised-clock data-acquisition problem; it is not
referenced further in this chapter.

%% ----------------------------------------------------------
\section{Class Imbalance Handling with SMOTE}
\label{sec:cl_imbalance}
%% ----------------------------------------------------------

\subsection{The Imbalance Problem}

As shown in Chapter~6, the activity-gated dataset is imbalanced:
\code{bruxism\_right} (325 epochs) is roughly $2\times$ more
frequent than \code{clinch} (165 epochs). Training a classifier
directly on this imbalanced data biases the model toward the
majority class, since minimising overall error is achieved by
favouring the class with the most examples.

\subsection{SMOTE}

The Synthetic Minority Over-sampling Technique
(\acrshort{smote})~\cite{chawla2002smote} addresses this imbalance
by generating synthetic minority-class samples rather than simply
duplicating existing ones. For each minority-class sample
$\mathbf{x}_i$, \acrshort{smote} selects one of its $k$ nearest
minority-class neighbours $\mathbf{x}_{nn}$ and generates a new
synthetic sample along the line segment connecting them:

\begin{equation}
\mathbf{x}_{\text{new}} = \mathbf{x}_i + \lambda
\left(\mathbf{x}_{nn} - \mathbf{x}_i\right),
\quad \lambda \sim \mathcal{U}(0, 1)
\label{eq:smote}
\end{equation}

\noindent where $\lambda$ is drawn uniformly at random from the
interval $[0, 1]$ and $k = 5$ neighbours are used by default. This
interpolation generates plausible new examples within the convex
hull of the existing minority-class samples, enriching the minority
class without the overfitting that naive duplication would cause.

\subsection{Application Protocol}

\acrshort{smote} is applied to the \textbf{training fold only},
after whichever split defines the current evaluation protocol
(GroupKFold or \acrshort{loso}). Applying \acrshort{smote} before
the split, or to the pooled dataset, would allow synthetic samples
derived from held-out examples to influence the training set ---
precisely the kind of preprocessing leak that
Section~\ref{sec:cl_leakage} audits for. By fitting \acrshort{smote}
strictly on the training fold, the held-out session remains a
genuine, untouched measure of generalisation performance. This
protocol is verified explicitly as part of the leakage audit
(Section~\ref{sec:cl_leakage}).

%% ----------------------------------------------------------
\section{The Evaluation Methodology Crisis}
\label{sec:cl_leakage}
%% ----------------------------------------------------------

This section documents what is, methodologically, the single most
important lesson of the project: an accuracy number is only as
trustworthy as the evaluation protocol that produced it, and the
naive protocol used in early development produced a number that
was not trustworthy at all.

\subsection{The Inflated Number}

Initial results using a random \code{StratifiedKFold} split gave an
apparent accuracy of approximately 80--85\%. As described in
Chapter~7, epochs are extracted with 75\% overlap between adjacent
windows (2.0\,s epochs, 0.5\,s step): two temporally adjacent
epochs share 300 of their 400 samples and are therefore
near-duplicates. If overlapping epochs are split randomly between
training and test folds, the classifier is effectively evaluated on
data it has already almost seen, and the reported accuracy is
substantially inflated by this \textbf{epoch-overlap leakage}
rather than reflecting genuine generalisation.

\subsection{The Leakage Audit}
\label{sec:cl_leakage_audit}

Rather than simply asserting that the naive protocol leaks, a
formal audit was run consisting of four independent diagnostic
tests, summarised in Table~\ref{tab:leakage_audit}.

\begin{table}[htbp]
\centering
\caption{The four-test leakage audit. Each test isolates a
different potential source of evaluation-time leakage.}
\label{tab:leakage_audit}
\renewcommand{\arraystretch}{1.4}
\begin{tabular}{p{3.4cm} p{3.6cm} p{5.0cm}}
\toprule
\textbf{Test} & \textbf{Purpose} & \textbf{Result} \\
\midrule
Strict cross-session (Session1+Session2 $\rightarrow$ Session3)
    & Tests true generalisation to an unseen recording day
    & 74.4\% ungated / 81\% gated --- does \emph{not} collapse
      $\rightarrow$ the underlying signal is real \\
Scaler / SelectKBest fit-on-train check
    & Tests for a preprocessing leak (statistics computed on
      data that includes the test fold)
    & \textbf{Clean} --- verified numerically; scaler and
      feature-selection statistics are always fit on the
      training fold only \\
Same-recording epochs present in both folds
    & Tests for the epoch-overlap leak described above
    & Headline (random) protocol: \textbf{yes, leaks}.
      Grouped protocols (GroupKFold, \acrshort{loso}):
      \textbf{no leak} \\
Shuffle test (class labels randomly permuted)
    & Tests for an implementation leak (any path by which the
      true label could reach the model other than through the
      features)
    & \textbf{21\% accuracy} ($\approx$ chance for 5 classes)
      $\rightarrow$ the pipeline does not manufacture signal \\
\bottomrule
\end{tabular}
\end{table}

The shuffle test deserves particular emphasis as a piece of
methodology worth generalising beyond this project: by permuting
the class labels while leaving every other part of the pipeline
untouched (same features, same scaler, same classifier, same
cross-validation splitting), any accuracy substantially above chance
would indicate that information about the true label is leaking
into the model through some channel other than genuine feature
content --- for example, through an index or ordering artefact.
Chance level for a balanced 5-class problem is 20\%; the observed
21\% confirms the pipeline is clean.

\textbf{Conclusion of the audit:} there is no implementation leak
and no preprocessing leak, but the headline (random-split) protocol
does leak via epoch overlap. This single finding is what
invalidates the $\sim$80--85\% figure obtained under naive
\code{StratifiedKFold} and motivates the grouped evaluation
protocols used for every result reported from this point in the
report onward.

\subsection{The Honest Protocols Going Forward}

Table~\ref{tab:honest_protocols} summarises the evaluation
protocols retained after the audit, and the role each plays.

\begin{table}[htbp]
\centering
\caption{Evaluation protocols used after the leakage audit. Random
StratifiedKFold is retained in this report only as a documented
negative example --- it must never be cited as a performance
number.}
\label{tab:honest_protocols}
\renewcommand{\arraystretch}{1.3}
\begin{tabular}{p{4.2cm} c c p{4.0cm}}
\toprule
\textbf{Protocol} & \textbf{Use} & \textbf{Leak-free?} &
\textbf{Property} \\
\midrule
Random StratifiedKFold & --- (documented only) & \ding{55} &
Inflated by epoch-overlap leakage; \textbf{do not cite} \\
GroupKFold (recording-level) & Development / ablation tuning &
\ding{51} & No overlap leakage, but not cross-day \\
\textbf{Leave-One-Session-Out} & \textbf{Final reporting} &
\ding{51} & Strictest: cross-day, deployment-realistic \\
\bottomrule
\end{tabular}
\end{table}

GroupKFold is used throughout development and ablation work in this
chapter because it is fast to iterate on and already eliminates the
epoch-overlap leak, at the cost of not testing across recording
days. \acrshort{loso} (Section~\ref{sec:cl_loso}) is strictly harder
--- and strictly more honest for this project's single-subject,
cross-day deployment scenario --- and is used for every number
that this report cites as a final result.

%% ----------------------------------------------------------
\section{Leave-One-Session-Out Evaluation Protocol}
\label{sec:cl_loso}
%% ----------------------------------------------------------

\subsection{Definition}

Leave-One-Session-Out (\acrshort{loso}) cross-validation treats
each of the three recording sessions (Chapter~6) as one fold. For
each session $s \in \{\text{Session1}, \text{Session2},
\text{Session3}\}$, the model is trained on the union of the
\emph{other two} sessions and evaluated on session $s$, which the
model has never seen in any form (not even through a shared scaler
or feature-selection statistic):

\begin{equation}
\text{Acc}_{\text{LOSO}} = \frac{1}{3}\sum_{s=1}^{3}
\text{Acc}\!\left(\mathcal{M}_{\setminus s},\, \mathcal{D}_s\right)
\label{eq:loso}
\end{equation}

\noindent where $\mathcal{M}_{\setminus s}$ is the model trained on
all sessions except $s$, and $\mathcal{D}_s$ is the held-out data
from session $s$. Because a session is an entire recording day ---
with its own electrode placement, impedance, and physiological
state --- \acrshort{loso} directly measures the quantity that
matters for deployment: how well a model calibrated on previous
days performs on a new day, without further recalibration.

\subsection{Why LOSO Is the Correct Protocol for This Project}

Since this is a \textbf{single-subject} design (Chapter~6), the
relevant notion of ``generalisation'' is not generalisation to a
new person, but generalisation to a new \emph{day} for the same
person --- which is exactly what \acrshort{loso} tests. This
distinguishes \acrshort{loso} from GroupKFold: GroupKFold holds out
individual \emph{recordings} (which may still come from the same
session as some training data), whereas \acrshort{loso} holds out
an entire \emph{session}, so nothing from that recording day ---
including its specific impedance and amplifier-gain conditions ---
is available to the model at evaluation time. This is the strictest
protocol available given three sessions, and the one most
representative of real deployment, where the user puts the headset
on for a new session and expects the previously trained model to
work with no additional calibration.

\subsection{Reference Numbers}

Table~\ref{tab:honest_numbers} consolidates the numbers referenced
throughout this chapter, so that any figure quoted later can be
traced back to the protocol that produced it.

\begin{table}[htbp]
\centering
\caption{Reference table for every accuracy figure used in this
report. Only the \acrshort{loso} row should ever be cited as
\emph{the} accuracy of the system.}
\label{tab:honest_numbers}
\renewcommand{\arraystretch}{1.3}
\begin{tabular}{p{4.5cm} p{4.7cm} c}
\toprule
\textbf{Protocol} & \textbf{What it measures} & \textbf{Number} \\
\midrule
Random StratifiedKFold & Inflated by overlap leakage &
$\sim$85\% --- \textbf{do not cite} \\
Shuffle test & Pipeline integrity check & 21\% ($\approx$chance)
--- clean \\
GroupKFold (recording-level) & Development metric (final hybrid) &
88.0\% \\
\textbf{LOSO (final, honest)} & \textbf{Cross-session deployment
accuracy} & \textbf{91.3\%} --- cite this \\
\bottomrule
\end{tabular}
\end{table}

%% ----------------------------------------------------------
\section{Baseline: Flat Random Forest Classifier}
\label{sec:cl_baseline}
%% ----------------------------------------------------------

Before any architectural changes, a conventional flat classifier
was trained on the full feature superset (Chapter~8) as a
reference point. \code{StandardScaler}, \code{SelectKBest} (mutual
information), and \acrshort{smote} were applied in sequence, ahead
of the classifier itself; Random Forest, \acrshort{svm} with an RBF
kernel, and \acrshort{lda} were all evaluated as candidate final
classifiers under this flat design, with Random Forest carried
forward as the production choice for the reasons summarised below.

\subsection{Why Random Forest}

\begin{itemize}[leftmargin=*]
    \item \textbf{Mixed feature handling:} Decision trees split on
    one feature at a time, so features of very different scales and
    distributions (bounded asymmetry indices, heavy-tailed
    kurtosis, strictly positive band powers) can be combined
    without normalisation issues.
    \item \textbf{Robustness to irrelevant features:} Random
    feature subsampling at each split means an irrelevant feature
    is simply not selected, which matters once the feature superset
    grows to 129 columns (Chapter~8) and different stages of the
    final architecture only need a subset of them.
    \item \textbf{Fast, deterministic inference:} prediction
    requires only traversing a fixed number of shallow trees,
    giving low, deterministic latency suitable for real-time
    operation (Chapter~12).
\end{itemize}

\subsection{Hyperparameters}

Table~\ref{tab:rf_hyperparams} lists the Random Forest
hyperparameters used for the Stage~1 router and, where applicable,
elsewhere in the pipeline.

\begin{table}[htbp]
\centering
\caption{Random Forest hyperparameters, used consistently as the
Stage~1 blink/muscle router throughout the hierarchical and hybrid
architectures.}
\label{tab:rf_hyperparams}
\renewcommand{\arraystretch}{1.3}
\begin{tabular}{ll}
\toprule
\textbf{Hyperparameter} & \textbf{Value} \\
\midrule
Number of trees (\code{n\_estimators}) & 300 \\
Max features per split & $\sqrt{93} \approx 10$ \\
Class weight & balanced \\
Bootstrap & True \\
Random seed & 42 (reproducibility) \\
\bottomrule
\end{tabular}
\end{table}

\subsection{Baseline Result}

Under \acrshort{loso}, the flat Random Forest baseline achieves
\textbf{85.1\% accuracy}, macro-F1 \textbf{0.837}. This is the
starting point against which every subsequent architectural change
in this chapter is measured (Section~\ref{sec:cl_progression}).

%% ----------------------------------------------------------
\section{The Hierarchical Architecture}
\label{sec:cl_hierarchical}
%% ----------------------------------------------------------

\subsection{Motivation: Features Cannot Be Applied Conditionally}

Chapter~8 documented a systematic set of extended feature
engineering experiments (\acrshort{dwt}, \acrshort{tkeo}, \acrshort{emg}
descriptors), all of which showed the same pattern: helpful for
either the blink classes or the muscle classes, never both, and
often actively harmful for the other half of the taxonomy. A single
flat classifier, trained on the concatenation of all features
across all classes, has no mechanism to apply a feature family
\emph{conditionally} --- it must treat every feature as globally
available or globally absent. This observation, not a search for
more features, is what motivated a hierarchical, physiology-routed
architecture.

\subsection{Design}

The hierarchical architecture routes each window through two
stages, illustrated in Figure~\ref{fig:hierarchical_arch}:

\begin{itemize}[leftmargin=*]
    \item \textbf{Stage 1 (blink vs.\ muscle):} a binary router
    trained on frontal/\acrshort{eog}-relevant features (Fp1, Fp2,
    Fz, F3, F4) that separates \{\code{single\_blink},
    \code{double\_blink}\} from \{\code{clinch},
    \code{bruxism\_left}, \code{bruxism\_right}\}.
    \item \textbf{Stage 2a (S vs.\ O):} conditioned on Stage~1
    routing to ``blink'', a sub-classifier trained \emph{only} on
    blink epochs, using blink peak features (Chapter~8).
    \item \textbf{Stage 2b (C vs.\ L vs.\ R):} conditioned on
    Stage~1 routing to ``muscle'', a sub-classifier trained
    \emph{only} on muscle epochs, using \acrshort{emg} descriptors
    and asymmetry features (Chapter~8).
\end{itemize}

\begin{figure}[htbp]
\centering
\begin{tikzpicture}[
    node distance=0.6cm and 1.2cm,
    box/.style={rectangle, rounded corners=3pt, draw=black,
                fill=blue!10, minimum width=4.2cm,
                minimum height=0.8cm, align=center, font=\small},
    arrow/.style={-{Stealth[length=5pt]}, thick}
]
\node[box, fill=green!15] (in) {129-Feature Superset (Chapter~8)};
\node[box, below=of in] (s1)
    {Stage 1: blink vs.\ muscle\\(frontal/EOG features)};
\node[box, fill=blue!10, below left=0.9cm and 0.4cm of s1] (s2a)
    {Stage 2a: S vs.\ O\\(blink peak features)};
\node[box, fill=orange!15, below right=0.9cm and 0.4cm of s1] (s2b)
    {Stage 2b: C vs.\ L vs.\ R\\(EMG descriptors + asymmetry)};
\node[box, fill=purple!10, below=2.4cm of s1, minimum width=6cm]
    (out) {Predicted class: O / S / C / L / R};

\draw[arrow] (in) -- (s1);
\draw[arrow] (s1.south) -| node[near start, above,
    font=\scriptsize]{blink} (s2a.north);
\draw[arrow] (s1.south) -| node[near start, above,
    font=\scriptsize]{muscle} (s2b.north);
\draw[arrow] (s2a.south) |- (out.west);
\draw[arrow] (s2b.south) |- (out.east);
\end{tikzpicture}
\caption{The hierarchical two-stage architecture. Each stage trains
on its own epoch subset and its own feature columns, so muscle
features can never reach the blink decision and vice versa.}
\label{fig:hierarchical_arch}
\end{figure}

Each stage trains only on its own epoch subset, so muscle features
can never reach the blink decision and vice versa. The full
129-feature superset is extracted once; each stage selects its
columns by name rather than the feature extraction being repeated
per stage.

\subsection{GroupKFold Development Results}

Table~\ref{tab:hierarchical_results} shows the per-stage and
end-to-end accuracy under GroupKFold, used for development and
ablation at this stage of the project.

\begin{table}[htbp]
\centering
\caption{Hierarchical architecture results under GroupKFold
(development protocol).}
\label{tab:hierarchical_results}
\renewcommand{\arraystretch}{1.3}
\begin{tabular}{l c}
\toprule
\textbf{Stage} & \textbf{Accuracy} \\
\midrule
Stage 1 (blink/muscle routing) & 98.7\% \\
Stage 2a (S vs.\ O)            & 76.4\% \\
Stage 2b (C vs.\ L vs.\ R)     & 87.8\% \\
\midrule
\textbf{End-to-end} & \textbf{82.2\%} (vs.\ 77.2\% flat) \\
\bottomrule
\end{tabular}
\end{table}

The biggest gains were in the blink classes (S~$+0.092$,
O~$+0.120$ F1), because Stage~1's near-perfect routing (98.7\%)
almost eliminates blink$\leftrightarrow$muscle confusion, and
Stage~2a sub-classifies blinks free of muscle-feature dilution.
Notably, the blink peak features (Chapter~8) --- which had failed
to help in the flat multi-class context due to label contamination
--- \emph{finally worked} in this pure 2-class S-vs-O setting,
because the contamination's confounding effect on decision
boundaries is far less severe when the only competing class is the
other blink type rather than five heterogeneous classes at once.

\subsection{LOSO Reconciliation}

On the strictest \acrshort{loso} protocol, the hierarchical model
initially \emph{regressed} by $-1.9$ points relative to the flat
baseline, despite its GroupKFold gains. An ablation traced this to
two session-fragile design choices:

\begin{itemize}[leftmargin=*]
    \item \textbf{\acrshort{tkeo} is the poison:} including the
    \acrshort{tkeo} feature (Chapter~8) in Stage~2b dragged
    \code{bruxism\_right} F1 down to 0.789 --- consistent with the
    Chapter~8 finding that \acrshort{tkeo}'s squared-amplitude
    measure does not survive cross-session drift. It was dropped.
    \item \textbf{\acrshort{emg} descriptors help cross-session:}
    despite hurting the flat model overall (Chapter~8), the
    \acrshort{emg} descriptor features improved clinch F1 from
    0.74 to 0.84 under \acrshort{loso} --- their complexity
    measures generalise well once isolated to the muscle-only
    Stage~2b. They were kept.
    \item \textbf{Stage 1 needed temporal channels:} adding T3/T4
    features to the Stage~1 router (in addition to the frontal
    features) took routing isolation from 98\% to 100\% under
    \acrshort{loso}, closing the last source of
    blink$\leftrightarrow$muscle cross-contamination.
\end{itemize}

After this fix, the hierarchical model matches the flat baseline at
\textbf{85.1\% \acrshort{loso}} accuracy, but with a better
macro-F1 (\textbf{0.839} vs.\ 0.837) and materially cleaner live
routing behaviour --- setting the stage for the two branch-specific
improvements described next.

%% ----------------------------------------------------------
\section{Riemannian Geometry for Muscle Classification}
\label{sec:cl_riemann}
%% ----------------------------------------------------------

\subsection{Motivation}

The \acrshort{tkeo} failure (Section~\ref{sec:cl_hierarchical};
Chapter~8) revealed that features built from \emph{absolute
amplitude} do not survive day-to-day amplitude drift between
recording sessions. Riemannian geometry offers an alternative: it
operates on the \textbf{inter-channel covariance structure} of an
epoch rather than on absolute amplitude, making it inherently more
robust to amplitude scaling --- and therefore a natural candidate
for the muscle branch (Stage~2b), where amplitude drift had already
been identified as the dominant cross-session failure mode.

\subsection{Mathematical Formulation}

For an epoch $\mathbf{E}_k \in \mathbb{R}^{C \times W}$ (in volts;
see the unit-correctness note below), the sample spatial covariance
matrix is:

\begin{equation}
\mathbf{C}_k = \frac{1}{W - 1} \mathbf{E}_k \mathbf{E}_k^{\top}
\in \mathbb{R}^{C \times C}
\label{eq:covariance}
\end{equation}

\noindent Sample covariance matrices estimated from short epochs
are noisy, so the Oracle Approximating Shrinkage
(\acrshort{oas})~\cite{barachant2012riemann} estimator is used
instead, which regularises the sample covariance toward a scaled
identity matrix:

\begin{equation}
\hat{\mathbf{C}}_k^{\text{OAS}} =
(1-\rho)\,\mathbf{C}_k + \rho\,\frac{\text{tr}(\mathbf{C}_k)}{C}\,
\mathbf{I}_C
\label{eq:oas}
\end{equation}

\noindent where $\rho \in [0,1]$ is a shrinkage coefficient chosen
automatically to minimise expected estimation error. Covariance
matrices are symmetric positive-definite (\acrshort{spd}) and do
not live on a flat Euclidean space; comparing them directly with
Euclidean distance is known to be sub-optimal. Instead, the
affine-invariant Riemannian distance between two \acrshort{spd}
matrices is:

\begin{equation}
\delta_R(\mathbf{C}_1, \mathbf{C}_2) =
\left\lVert \log\!\left(\mathbf{C}_1^{-1/2}\,\mathbf{C}_2\,
\mathbf{C}_1^{-1/2}\right)\right\rVert_F
\label{eq:riemann_dist}
\end{equation}

To use standard Euclidean classifiers (such as \acrshort{svc}),
each covariance matrix is projected onto the \textbf{tangent
space} at a reference point $\bar{\mathbf{C}}$ (typically the
geometric mean of the training covariances):

\begin{equation}
\mathbf{s}_k = \text{upper}\!\left(
\bar{\mathbf{C}}^{1/2}\,
\log\!\left(\bar{\mathbf{C}}^{-1/2}\,\mathbf{C}_k\,
\bar{\mathbf{C}}^{-1/2}\right)\,
\bar{\mathbf{C}}^{1/2}\right)
\label{eq:tangent_space}
\end{equation}

\noindent where $\text{upper}(\cdot)$ vectorises the upper
triangular part of the resulting symmetric matrix. The tangent
vectors $\mathbf{s}_k \in \mathbb{R}^{C(C+1)/2}$ live in a flat
Euclidean space and can be classified directly with any standard
classifier.

\subsection{Method and Implementation}

The muscle branch (Stage~2b) pipeline is:

\begin{lstlisting}[language=Python,
    caption={Riemannian muscle sub-classifier pipeline
    (\code{riemann.py}).},
    label={lst:riemann}]
from pyriemann.estimation import Covariances
from pyriemann.tangentspace import TangentSpace
from sklearn.svm import SVC
from sklearn.pipeline import make_pipeline

muscle_pipeline = make_pipeline(
    Covariances(estimator='oas'),
    TangentSpace(metric='riemann'),
    SVC(kernel='linear', C=1.0)
)
\end{lstlisting}

Unlike the hand-crafted feature branches, this pipeline operates on
\textbf{raw epochs} directly (covariance is computed across
channels from the unfiltered-by-feature-extraction signal), not on
engineered feature columns. It is applied only to epochs already
routed to ``muscle'' by Stage~1.

\subsection{Results}

Table~\ref{tab:riemann_results} compares the hierarchical
(hand-crafted feature) muscle branch against the Riemannian muscle
branch, both under \acrshort{loso}.

\begin{table}[htbp]
\centering
\caption{Riemannian muscle branch vs.\ hand-crafted-feature muscle
branch, LOSO protocol. Per-class columns are F1 scores.}
\label{tab:riemann_results}
\renewcommand{\arraystretch}{1.3}
\begin{tabular}{l c c c c}
\toprule
\textbf{Model} & \textbf{Overall} & \textbf{C} & \textbf{L} &
\textbf{R} \\
\midrule
Hierarchical (hand-crafted features) & 85.1\% & 0.83 & 0.90 & 0.94 \\
\textbf{Riemannian hybrid} & \textbf{89.2\%} & \textbf{0.952} &
\textbf{0.984} & \textbf{0.983} \\
\bottomrule
\end{tabular}
\end{table}

The entire $+4.1$-point gain (85.1\%~$\rightarrow$~89.2\%) comes
from the muscle classes, confirming the central hypothesis: spatial
covariance structure survives cross-session amplitude drift far
better than any amplitude-based feature tried in Chapter~8.

\subsection{Key Negative Findings}

\begin{itemize}[leftmargin=*]
    \item \textbf{Pure Riemannian is not viable end-to-end:}
    applying the Riemannian pipeline to \emph{all five} classes
    (not just the muscle branch) is near-perfect on muscle classes
    but destroys the blink classes, because covariance structure is
    blind to the temporal single-vs-double blink pattern that
    distinguishes them --- covariance summarises spatial
    relationships, not event timing. The hybrid (branch-specific)
    architecture is therefore essential, not optional.
    \item \textbf{\code{tsupdate=True} is a trap:} enabling online
    tangent-space recentring gives the best score under
    \acrshort{loso} but \emph{collapses} under GroupKFold (clinch
    F1 drops to 0.653), indicating it overfits to session-specific
    structure rather than genuinely improving the representation.
    Plain \code{TangentSpace} without online updating is robust on
    both protocols and is what is used in production.
    \item \textbf{Unit correctness:} the tangent space is fit in
    volts, but the raw \acrshort{edf} signal and upstream pipeline
    (Chapter~7) work in microvolts throughout. Inference must
    explicitly convert $\mu\text{V} \rightarrow \text{V}$ before
    computing covariance; a silent unit mismatch here would not
    raise an error but would silently wreck muscle-class
    predictions, since the covariance scale would be off by a
    factor of $10^{12}$.
\end{itemize}

%% ----------------------------------------------------------
\section{The Blink Ceiling: Five Attempts}
\label{sec:cl_dtw}
%% ----------------------------------------------------------

While the muscle branch improved substantially, the blink classes
(S/O) remained stuck at an F1 ceiling of roughly 0.74/0.78
throughout the hierarchical and Riemannian-hybrid stages. Five
independent methods were attempted for single-vs-double blink
discrimination; four failed, and the diagnostic \emph{reason} they
failed turned out to be the more important finding.
Table~\ref{tab:blink_attempts} summarises all five.

\begin{table}[htbp]
\centering
\caption{The five attempts at blink sub-classification. Only the
fifth, a fundamentally different approach, succeeded.}
\label{tab:blink_attempts}
\renewcommand{\arraystretch}{1.3}
\begin{tabular}{c p{3.2cm} c p{5.0cm}}
\toprule
\textbf{\#} & \textbf{Method} & \textbf{Result} &
\textbf{Why it failed} \\
\midrule
1 & Raw peak-count & Failed & Counts compressed together;
    single $\approx$ double \\
2 & $\mu$V peak-count & Failed & Z-scoring (Chapter~7) breaks the
    absolute amplitude threshold the method relies on \\
3 & Sliding-window count (moving std.) & Failed (70.2\% vs.\
    77.8\%) & 2\,s windows are not event-locked to individual
    blinks \\
4 & Event-locked epoching & Failed at validation gate & Diagnostic
    revealed single-blink recordings are contaminated
    (below) \\
5 & \textbf{DTW shape-matching} & \textbf{Success ($+2.1\%$)} &
    Shape survives contamination that counting cannot \\
\bottomrule
\end{tabular}
\end{table}

\subsection{The Diagnostic Discovery (Attempts 1--4)}

Attempt 4's validation gate --- which checks how many blinks an
automated detector actually finds inside each nominally-labelled
recording --- produced the finding already introduced in Chapter~6:

\begin{table}[htbp]
\centering
\caption{Blinks detected per nominal gesture trial, by class. This
diagnostic, run while building Attempt 4, is what revealed the
single-blink contamination problem (Chapter~6).}
\label{tab:blink_diagnostic}
\renewcommand{\arraystretch}{1.3}
\begin{tabular}{l c c p{4.5cm}}
\toprule
\textbf{Class} & \textbf{Expected} & \textbf{Detected (mean)} &
\textbf{Quality} \\
\midrule
\code{single\_blink} & 1 & 2.60 & Only 39\% clean; 45\% are
    $\geq$1\,s clusters of $\geq$3 blinks \\
\code{double\_blink} & 2 & 2.07 & 55\% correct \ding{51} \\
\midrule
\multicolumn{4}{l}{\footnotesize (\code{triple\_blink}, detected
    at 3.08, was clean but is not part of the final taxonomy ---
    Chapter~6)} \\
\bottomrule
\end{tabular}
\end{table}

The detector correctly counts blinks in the \code{double\_blink}
recordings; the problem is that a substantial fraction of the
\code{single\_blink} recordings are \emph{physically} contaminated
with multi-blink bursts (Chapter~6). \textbf{When one class's
recordings contain genuine instances of another class, no counting
method --- however carefully engineered --- can separate them.}
This conclusion is not a hypothesis; it is proven empirically across
four independently-designed counting approaches, all of which fail
for the same underlying reason.

\subsection{The Breakthrough: DTW Shape-Matching (Attempt 5)}

The insight behind Attempt 5 is to stop counting discrete events and
instead compare \textbf{waveform shape} directly. Dynamic Time
Warping (\acrshort{dtw})~\cite{sakoe1978dtw, berndt1994dtw} computes
an elastic, non-linear alignment between two time series that
minimises cumulative pointwise distance:

\begin{equation}
D(i,j) = d(x_i, y_j) +
\min\!\big(D(i{-}1,j),\, D(i,j{-}1),\, D(i{-}1,j{-}1)\big)
\label{eq:dtw}
\end{equation}

\noindent with $D(0,0) = 0$ and $d(x_i,y_j) = (x_i - y_j)^2$ the
local pointwise cost. A Sakoe--Chiba band of radius $r$ constrains
the alignment path to $|i-j| \leq r$, preventing pathological
warps and bounding computation to $O(rW)$ per pair instead of
$O(W^2)$. Under this metric, a single blink (one bump), a
double blink (two bumps), and a genuinely-clean or contaminated
window each have a distinct \emph{shape}, largely independent of
exactly when each bump occurs --- so a contaminated single-blink
window, while it may fool a peak counter, still has a shape closest
to other single-blink exemplars rather than to double-blink ones.

\textbf{Method:}

\begin{lstlisting}[language=Python,
    caption={DTW blink sub-classifier pipeline (Stage~2a,
    production).},
    label={lst:dtw}]
from tslearn.preprocessing import TimeSeriesScalerMeanVariance
from tslearn.neighbors import KNeighborsTimeSeriesClassifier

blink_pipeline_steps = [
    TimeSeriesScalerMeanVariance(),
    KNeighborsTimeSeriesClassifier(
        n_neighbors=5,
        metric="dtw",
        metric_params={"sakoe_chiba_radius": 10}
    )
]
\end{lstlisting}

\noindent applied to the z-scored frontal waveform (the mean of
Fp1, Fp2, Fz) of each epoch already routed to ``blink'' by Stage~1.

\textbf{Critical sub-finding --- exemplar voting beats averaging:}
\acrshort{dtw} \emph{barycenter} templates (one averaged prototype
per class, compared by \acrshort{dtw} distance) were also tried and
\emph{failed} outright (50.5\% accuracy), because the contaminated
\code{single\_blink} class has no meaningful average shape to
compute --- averaging a clean single-bump trial with a contaminated
triple-bump trial produces a template that resembles neither.
\textbf{$k$-nearest-neighbour exemplar voting} is what survives the
contamination: a noisy single-blink window still lands nearest to
\emph{some} genuinely clean single-blink exemplar in the training
set, even when it cannot be meaningfully averaged or counted.

\subsection{Results}

Table~\ref{tab:dtw_results} shows the final blink-branch comparison
under \acrshort{loso}.

\begin{table}[htbp]
\centering
\caption{DTW blink branch vs.\ Riemannian-hybrid blink branch
(which still used hand-crafted peak features for Stage~2a), LOSO
protocol.}
\label{tab:dtw_results}
\renewcommand{\arraystretch}{1.3}
\begin{tabular}{l c c c c}
\toprule
\textbf{Model} & \textbf{Overall} & \textbf{macro-F1} &
\textbf{S (double)} & \textbf{O (single)} \\
\midrule
Riemann hybrid & 89.2\% & 0.888 & 0.742 & 0.778 \\
\textbf{DTW hybrid (final)} & \textbf{91.3\%} & \textbf{0.913} &
\textbf{0.806} & \textbf{0.841} \\
\bottomrule
\end{tabular}
\end{table}

The $+2.1$-point gain (89.2\%~$\rightarrow$~91.3\%) is entirely
attributable to the blink branch, and completes the final
production architecture.

%% ----------------------------------------------------------
\section{Final Architecture: The Three-Way Hybrid}
\label{sec:cl_architecture}
%% ----------------------------------------------------------

The production model, serialised as \code{dtw\_hybrid\_model.pkl},
routes each gesture family to the classification method
physiologically matched to it, illustrated in
Figure~\ref{fig:final_arch}.

\begin{figure}[htbp]
\centering
\begin{tikzpicture}[
    node distance=0.55cm and 1.0cm,
    box/.style={rectangle, rounded corners=3pt, draw=black,
                fill=blue!10, minimum width=4.6cm,
                minimum height=0.8cm, align=center, font=\small},
    arrow/.style={-{Stealth[length=5pt]}, thick}
]
\node[box, fill=green!15] (win) {Window (2.0\,s)};
\node[box, below=of win] (gate) {Activity Gate};
\node[box, below=of gate] (s1)
    {Stage 1: blink vs.\ muscle (features)};
\node[box, fill=blue!10, below left=0.9cm and 0.3cm of s1] (s2a)
    {Stage 2a: S vs.\ O\\\acrshort{dtw}-KNN on waveform\\(shape
    $>$ count)};
\node[box, fill=orange!15, below right=0.9cm and 0.3cm of s1] (s2b)
    {Stage 2b: C vs.\ L vs.\ R\\Riemannian covariance\\(covariance
    $>$ amplitude)};
\node[box, fill=purple!10, below=2.6cm of s1, minimum width=6.4cm]
    (merge) {Confidence gate ($>0.60$)};
\node[box, fill=purple!10, below=of merge, minimum width=6.4cm]
    (fsm) {\acrshort{fsm} (3 consecutive identical predictions)};
\node[box, fill=red!10, below=of fsm, minimum width=6.4cm]
    (out) {\acrshort{uart} output: O / S / C / L / R (Chapter~12)};

\draw[arrow] (win) -- (gate);
\draw[arrow] (gate) -- (s1);
\draw[arrow] (s1.south) -| node[near start, above,
    font=\scriptsize]{blink} (s2a.north);
\draw[arrow] (s1.south) -| node[near start, above,
    font=\scriptsize]{muscle} (s2b.north);
\draw[arrow] (s2a.south) |- (merge.west);
\draw[arrow] (s2b.south) |- (merge.east);
\draw[arrow] (merge) -- (fsm);
\draw[arrow] (fsm) -- (out);
\end{tikzpicture}
\caption{The final three-way hybrid architecture. Training-time
routing (Stage 1, 2a, 2b) is the focus of this chapter; the
confidence gate, FSM, and UART output are real-time deployment
mechanics covered in full in Chapter~12.}
\label{fig:final_arch}
\end{figure}

Table~\ref{tab:why_each_method} summarises why each branch uses the
method it does, tying together the narrative of this chapter.

\begin{table}[htbp]
\centering
\caption{Why each branch of the final architecture uses the method
it does.}
\label{tab:why_each_method}
\renewcommand{\arraystretch}{1.4}
\begin{tabular}{p{2.3cm} p{3.4cm} p{6.0cm}}
\toprule
\textbf{Branch} & \textbf{Method} &
\textbf{Physiological rationale} \\
\midrule
Stage 1 & Hand-crafted features (Chapter~8) & \acrshort{eog} vs.\
    \acrshort{emg} separation is an easy problem (98--100\%
    accuracy); no need for a more elaborate representation. \\
Stage 2a (blink) & \acrshort{dtw} shape-matching & Temporal
    pattern (bump count and shape) survives label contamination
    that defeats every counting-based method. \\
Stage 2b (muscle) & Riemannian covariance & Spatial covariance
    structure survives cross-session amplitude drift that defeats
    every amplitude-based feature. \\
\bottomrule
\end{tabular}
\end{table}

%% ----------------------------------------------------------
\section{Performance Progression}
\label{sec:cl_progression}
%% ----------------------------------------------------------

Table~\ref{tab:progression} traces the cumulative effect of every
architectural decision documented in this chapter, from the flat
baseline to the final production model, all under the honest
\acrshort{loso} protocol.

\begin{table}[htbp]
\centering
\caption{Cumulative performance progression under LOSO. Each row
adds exactly one architectural change over the row above it.}
\label{tab:progression}
\renewcommand{\arraystretch}{1.3}
\begin{tabular}{l c c c}
\toprule
\textbf{Architecture} & \textbf{Overall} & \textbf{macro-F1} &
\textbf{$\Delta$ vs.\ previous} \\
\midrule
Flat Random Forest & 85.1\% & 0.837 & --- \\
Hierarchical (physiology-routed) & 85.1\% & 0.839 &
    $+0.0$pt (better macro-F1, cleaner routing) \\
+ Riemannian muscle branch & 89.2\% & 0.888 & $+4.1$pt \\
\textbf{+ DTW blink branch (final)} & \textbf{91.3\%} &
    \textbf{0.913} & \textbf{$+2.1$pt} \\
\bottomrule
\end{tabular}
\end{table}

Table~\ref{tab:final_f1} gives the final per-class F1 scores of the
production model; the full confusion matrix and per-class
precision/recall breakdown are presented in Chapter~10.

\begin{table}[htbp]
\centering
\caption{Final per-class F1 scores, DTW hybrid model, LOSO
protocol.}
\label{tab:final_f1}
\renewcommand{\arraystretch}{1.3}
\begin{tabular}{c l c}
\toprule
\textbf{Command} & \textbf{Class} & \textbf{F1} \\
\midrule
S & \code{double\_blink}  & 0.806 \\
O & \code{single\_blink}  & 0.841 \\
C & \code{clinch}         & 0.952 \\
L & \code{bruxism\_left}  & 0.984 \\
R & \code{bruxism\_right} & 0.983 \\
\bottomrule
\end{tabular}
\end{table}

The pattern across the whole progression is consistent with the
chapter's central argument: every gain came from matching a
classification method to the physiology of the signal it processes
(Riemannian covariance for amplitude-drift-prone muscle signals,
\acrshort{dtw} shape-matching for contamination-prone blink
signals), not from adding more features or more classifier
complexity to a single flat model.

%% ----------------------------------------------------------
\section{Chapter Summary}
\label{sec:cl_summary}
%% ----------------------------------------------------------

This chapter presented the complete classification and training
methodology. An evaluation methodology crisis was identified and
resolved: naive random cross-validation inflated accuracy by
exploiting 75\% epoch overlap, diagnosed through a four-test
leakage audit (strict cross-session check, preprocessing-leak
check, overlap check, and a shuffle test at 21\% $\approx$ chance),
leading to the adoption of Leave-One-Session-Out as the sole
protocol for honest reporting. A flat Random Forest baseline
(85.1\% \acrshort{loso}) was superseded by a hierarchical,
physiology-routed architecture (85.1\%, improved macro-F1 0.839),
motivated by Chapter~8's finding that no single feature family
helps both blink and muscle classes simultaneously. The muscle
branch was then upgraded to Riemannian covariance geometry
(+4.1\,pt, 89.2\%), exploiting covariance structure's robustness
to cross-session amplitude drift. Finally, after four
counting-based methods failed to resolve the blink sub-classification
ceiling --- failures traced to a critical data-quality finding
(Chapter~6) that $\sim$45\% of \code{single\_blink} recordings are
contaminated with multi-blink bursts --- Dynamic Time Warping
shape-matching with $k$-nearest-neighbour exemplar voting resolved
it (+2.1\,pt, \textbf{91.3\% final \acrshort{loso} accuracy,
macro-F1 0.913}). The resulting three-way hybrid architecture
routes each gesture family to the method physiologically matched to
it: hand-crafted features for the easy blink/muscle split,
\acrshort{dtw} for contamination-resistant blink shape-matching,
and Riemannian geometry for drift-resistant muscle covariance. The
following chapter presents the detailed results and analysis of
this final model, including the full confusion matrix and
per-class precision/recall breakdown.